% ID: MAT-TRI-FUN-001
% Titolo: Modello periodico di un silo
% Disciplina: Matematica
% Area: Trigonometria
% Argomento: Funzioni goniometriche
% Classe: Quarta secondaria di secondo grado
% Difficolta: 3
% Tipo: modellizzazione
% Tag: matematica, trigonometria, funzioni_goniometriche
% Fonte: verifica_funzioni_goniometriche_anonimizzata
% Provenienza: materiale fornito dal docente; testo anonimizzato; soluzione redatta per Exergo

\begin{esercizio}
Un silo per lo stoccaggio delle granaglie viene riempito e svuotato periodicamente durante l'anno. La relazione che lega la massa $m$ (in tonnellate) e il tempo $t$ (in mesi) pu\`o essere approssimata dalla funzione
    \[
        m(t)=20+15\sin\left(\frac{\pi}{3}t\right).
    \]
    \begin{enumerate}
        \item Il silo viene svuotato completamente? In caso di risposta negativa, qual \`e la scorta di mangime a disposizione?
        \item Quante volte viene riempito il silo in un anno?
        \item Disegna il grafico di $m(t)$ corrispondente a un anno.
    \end{enumerate}
\end{esercizio}

\begin{soluzione}
L'oscillazione è compresa fra $20-15=5$ e $20+15=35$ tonnellate:
il silo non si svuota e la scorta minima è 5 tonnellate.
Il periodo è $2\pi/(\pi/3)=6$ mesi, quindi in un anno si hanno due cicli e due massimi.
Per tracciare il grafico su $[0,12]$, usa i punti
$(0,20),(1{,}5,35),(3,20),(4{,}5,5),(6,20),(7{,}5,35),(9,20),(10{,}5,5),(12,20)$.
Il modello continuo non descrive singole operazioni discrete di carico.
\end{soluzione}
